\documentclass[11pt]{article}
\usepackage[margin=2.2cm]{geometry}
\usepackage{amsmath,amssymb}
\usepackage{booktabs}
\usepackage{parskip}
\usepackage{fancyvrb}
\usepackage{temperleylieb}
\usepackage[hidelinks]{hyperref}

\DefineVerbatimEnvironment{code}{Verbatim}
  {frame=leftline, framerule=1pt, rulecolor=\color{black!25},
   xleftmargin=1em, fontsize=\small, formatcom=\color{black!80}}
\newenvironment{result}{\par\smallskip\noindent\hspace*{1em}}{\par\medskip}
\newcommand\opt[1]{\texttt{#1}}
\newcommand\cs[1]{\texttt{\textbackslash#1}}
\newcommand\meta[1]{\ensuremath{\langle}\textit{#1}\ensuremath{\rangle}}
\newcommand\marg[1]{\texttt{\{}\meta{#1}\texttt{\}}}
\newcommand\oarg[1]{\texttt{[}\meta{#1}\texttt{]}}

\title{The \textsf{temperleylieb} package\\[2pt]
  \large Temperley--Lieb diagrams in \LaTeX{} with Ti\emph{k}Z}
\date{Version 1.0, 2026/10/21}
\author{}

\begin{document}
\maketitle

\noindent
The Temperley--Lieb algebra $TL_n(\delta)$ is generated by the identity
$1 = e_0$ and $e_1,\dots,e_{n-1}$ subject to the Jones relations
$e_i^2 = \delta e_i$, $e_ie_{i\pm1}e_i = e_i$ and $e_ie_j = e_je_i$ for
$|i-j|>1$. Elements are drawn as diagrams of $n$ non-crossing strands; the
product is concatenation, and every closed loop is replaced by a factor
$\delta$. This package draws such diagrams from words typed exactly as you
would write them in mathematics:
\[
  \TL{5}{e_1 e_4 e_3 e_2} \times \TL{5}{e_2 e_4 e_3}
  = \TL[boundaries]{5}{e_1 e_4 e_3 e_2 e_2 e_4 e_3}
  = \TL[reduce]{5}{e_1 e_4 e_3 e_2 e_2 e_4 e_3}
\]

\section{Quick start}

\begin{code}
\usepackage{temperleylieb}   % loads TikZ
...
\TL{5}{e_1 e_4 e_3 e_2}
\end{code}
\begin{result}
  \TL{5}{e_1 e_4 e_3 e_2}
\end{result}

The first mandatory argument is the number $n$ of strands, the second one is
the word. A word is a sequence of generators \verb|e_i| (or \verb|e_{12}| for
two-digit indices); \verb|e_0| and \verb|1| both denote the identity. Spaces,
\verb|*|, \verb|.|, \verb|\cdot|, \verb|\times| and \verb|\circ| between
generators are ignored, so \verb|e_1\cdot e_2| works as well.

\section{The three features}

\subsection{Composition (multiplication in the algebra)}

By default the generators of the word are \emph{composed}: they are glued
one after the other into a single diagram, and loops created along the way are
drawn. Add \opt{reduce} to get the reduced diagram (the basis element) with
the power of $\delta$ in front; \opt{boundaries} shows where the generators
were glued.

\begin{code}
\( \TL[boundaries]{3}{e_1 e_1 e_2 e_1} = \TL[reduce]{3}{e_1 e_1 e_2 e_1} \)
\end{code}
\begin{result}
  \( \TL[boundaries]{3}{e_1 e_1 e_2 e_1} = \TL[reduce]{3}{e_1 e_1 e_2 e_1} \)
\end{result}

After every diagram, \cs{TLlastloops} expands to the number of closed loops
that were found (here \TLlastloops).

\subsection{Displaying generators side by side without composing}

With \opt{separate}, each generator becomes its own diagram, laid out in a
row. Braces group generators that should still be composed into one
factor. Alternatively, in the default \opt{compose} mode, a bar \verb+|+
splits the word into factors. The key \opt{separator} puts a symbol between
the factors, and \opt{names} writes the generators underneath.

\begin{code}
\TLgenerators[names]{5}
\TL[separate, names]{4}{e_1 e_2 {e_1 e_3} 1}
\TL[separator=\times]{5}{e_1 e_4 e_3 e_2 | e_2 e_4 e_3}
\end{code}
\begin{result}
  \TLgenerators[names]{5}\qquad
  \TL[separate, names]{4}{e_1 e_2 {e_1 e_3} 1}\qquad
  \TL[separator=\times]{5}{e_1 e_4 e_3 e_2 | e_2 e_4 e_3}
\end{result}

\cs{TLgenerators}\oarg{options}\marg{n} draws $1, e_1, \dots, e_{n-1}$;
the starred form omits the identity. Combining \opt{separate} with
\opt{reduce} reduces each factor on its own, which is handy for lists of basis
elements:

\begin{code}
\TL[separate, reduce, names, separator={,}]{3}{1 e_1 e_2 {e_1 e_2} {e_2 e_1}}
\end{code}
\begin{result}
  \TL[separate, reduce, names, separator={,}]{3}{1 e_1 e_2 {e_1 e_2} {e_2 e_1}}
\end{result}

\subsection{Horizontal or vertical}

Diagrams are \opt{horizontal} by default (composition runs from left to
right, strand~1 on top). With \opt{vertical}, composition runs from top to
bottom and strand~1 is on the left. \opt{reverse} flips the direction of
composition (right to left, or bottom to top).

\begin{code}
\TL[vertical, labels]{4}{e_1 e_2 e_3}  \TL[vertical, reduce]{4}{e_1 e_2 e_3}
\TL[vertical, reverse]{4}{e_1 e_2 e_3}  \TL[reverse, names]{4}{e_1 e_2 e_3}
\TL[vertical, separate, names]{4}{e_1 e_2 e_3}
\end{code}
\begin{result}
  \TL[vertical, labels]{4}{e_1 e_2 e_3}\quad
  \TL[vertical, reduce]{4}{e_1 e_2 e_3}\quad
  \TL[vertical, reverse]{4}{e_1 e_2 e_3}\quad
  \TL[reverse, names]{4}{e_1 e_2 e_3}\quad
  \TL[vertical, separate, names]{4}{e_1 e_2 e_3}
\end{result}

\section{Reference}

\subsection{Commands}

\begin{description}
  \item[\cs{TL}\oarg{options}\marg{n}\marg{word}] Draws a diagram in
    $TL_n$. Works in text and in math; the diagram is vertically centred on
    the math axis. \cs{TLdiagram} is the same command under a longer name
    (\cs{TL} is only defined if no other package has taken it).
  \item[\cs{TLgenerators}\texttt{*}\oarg{options}\marg{n}] Draws
    $1,e_1,\dots,e_{n-1}$ side by side (without $1$ in the starred form).
  \item[\cs{tlset}\marg{options}] Sets options for all following diagrams in
    the current group.
  \item[\cs{TLlastloops}] Number of closed loops in the last diagram.
\end{description}

\subsection{Options}

\begin{center}
\small
\begin{tabular}{@{}lll@{}}
  \toprule
  Option & Default & Meaning \\
  \midrule
  \opt{compose} / \opt{separate} & \opt{compose} & glue the word, or draw each generator alone \\
  \opt{reduce} & false & draw reduced diagram(s) with a power of $\delta$ \\
  \opt{horizontal} / \opt{vertical} & \opt{horizontal} & direction of composition \\
  \opt{reverse} & false & compose right-to-left / bottom-to-top \\
  \opt{names} & false & write the generators below the diagram \\
  \opt{labels} & false & number the strands \\
  \opt{boundaries} & false & dashed lines between composed generators \\
  \opt{separator=}\meta{math} & empty & symbol between separate factors \\
  \opt{coefficient} & true & print $\delta^k$ in front of reduced diagrams \\
  \opt{loop value=}\meta{math} & \verb|\delta| & symbol used for a loop \\
  \opt{box=frame|fill|none} & \opt{frame} & background of each diagram \\
  \opt{strand sep=}\meta{number} & 0.4 & distance between strands (cm) \\
  \opt{length=}\meta{number} & 0.6 & length of one generator (cm) \\
  \opt{gap=}\meta{number} & 0.25 & space between separate factors (cm) \\
  \opt{scale=}\meta{number} & 1 & scales the picture \\
  \opt{baseline=}\meta{dim} & \verb|-0.5ex| & TikZ baseline (picture is centred at $y=0$) \\
  \opt{strand style=}\meta{tikz} & & extra TikZ options for strands \\
  \opt{box style=}\meta{tikz} & & extra TikZ options for the boxes \\
  \opt{tikz=}\meta{tikz} & & extra options for the \texttt{tikzpicture} \\
  \bottomrule
\end{tabular}
\end{center}

\subsection{Styles}

The look is controlled by TikZ styles that you can change with
\cs{tikzset}: \texttt{tl strand}, \texttt{tl box}, \texttt{tl box frame},
\texttt{tl box fill}, \texttt{tl boundary}, \texttt{tl label},
\texttt{tl name}, \texttt{tl separator} and \texttt{tl coefficient}. The
colour \texttt{tllightblue} is used by \opt{box=fill}.

\begin{code}
\tlset{box=fill, strand style={line width=1.4pt}, strand sep=0.35}
\TL[reduce, loop value=(q+q^{-1})]{4}{e_2 e_1 e_3 e_2 e_2}
\TL[box=none, strand style={red}, scale=0.7]{4}{e_2 e_1 e_3 e_2 e_2}
\end{code}
\begin{result}
  \begingroup
  \tlset{box=fill, strand style={line width=1.4pt}, strand sep=0.35}
  \TL[reduce, loop value=(q+q^{-1})]{4}{e_2 e_1 e_3 e_2 e_2}\quad
  \TL[box=none, strand style={red}, scale=0.7]{4}{e_2 e_1 e_3 e_2 e_2}
  \endgroup
\end{result}

\section{The Jones relations}

\begin{code}
\tlset{box=fill}
\[ \TL{3}{e_1 e_1} = \delta \TL{3}{e_1}, \qquad
   \TL{3}{e_1 e_2 e_1} = \TL{3}{e_1}, \qquad
   \TL{4}{e_1 e_3} = \TL{4}{e_3 e_1} \]
\end{code}
\begingroup
\tlset{box=fill}
\[ \TL{3}{e_1 e_1} = \delta \TL{3}{e_1}, \qquad
   \TL{3}{e_1 e_2 e_1} = \TL{3}{e_1}, \qquad
   \TL{4}{e_1 e_3} = \TL{4}{e_3 e_1} \]
\endgroup

\section{How the reduction works}

A word $e_{g_1}\cdots e_{g_m}$ is turned into a graph on the points
$(c, j)$, $0\le c\le m$, $1\le j\le n$: generator $e_{g_k}$ joins column
$k-1$ to column $k$ by straight strands, except for strands $g_k$ and
$g_k+1$, which are joined by a cup on column $k-1$ and a cap on column $k$.
Following each strand from the boundary gives the planar matching of the
reduced diagram; the remaining (internal) components are the closed loops.
Reduced cups and caps are drawn with a depth proportional to their nesting
level, so that nested arcs never touch.

\end{document}
